\section*{Chapter \ref{ch:g12:project-route-planner} --- Project: A Route Planner}

\begin{solution}{exo:g12:project-route-planner:1}
Direct: $3.6 \times 1730 / 30 = 207.6$, so $208$ s. Through the Theatre: $3.6 \times 890 / 30 = 106.8$, so $107$ s,
then $3.6 \times 1540 / 70 = 79.2$, so $79$ s: $186$ s in all. The call returns
\texttt{([\textquotesingle{}Bakery Corner\textquotesingle{}, \textquotesingle{}Theatre\textquotesingle{}, \textquotesingle{}South School\textquotesingle{}], 186)}: the longer route is faster.
\end{solution}

\begin{solution}{exo:g12:project-route-planner:2}
$3.6 \times 1830 / 30 = 219.6$, rounded to $220$ seconds, more than three and a half minutes for under two kilometres.
\end{solution}

\begin{solution}{exo:g12:project-route-planner:3}
\code{(20, 1830)}. The count is in the output of milestone~1: $20$ roads at $30$ km/h. The longest road of the town,
$1830$ m, is at $30$ km/h (the query of the town's tables), so it is also the longest of those $20$.
\end{solution}

\begin{solution}{exo:g12:project-route-planner:4}
Two arcs for each of the $52$ two-way roads and one for the one-way road: $105$. With $k$ one-way roads,
$2(53 - k) + k = 106 - k$.
\end{solution}

\begin{solution}{exo:g12:project-route-planner:5}
\smallskip\noindent\hspace*{2em}\code{def place_id(conn, name):}

\noindent\hspace*{4em}\code{row = conn.execute("SELECT place_id FROM place WHERE name = ?",}

\noindent\hspace*{12em}\code{(name,)).fetchone()}

\noindent\hspace*{4em}\code{if row is None:}

\noindent\hspace*{6em}\code{raise ValueError(f"no place called {name!r}")}

\noindent\hspace*{4em}\code{return row[0]}

\smallskip\noindent The name goes in as a parameter, never pasted into the text of the query, against \omterm{def:g12:relational-databases-sql:injection}{SQL injection}.
\code{fetchone} returns \code{None} when no row matches. The file is \code{exercises/rp_ex.py}.
\end{solution}

\begin{solution}{exo:g12:project-route-planner:6}
The \omterm{def:g10:graphs-first-look:graph}{graph} now has $106$ arcs, and \code{one_way_arcs()} returns \code{[]} instead of \code{[(1, 3)]}: both checks of
milestone~2 that concern road~$1$ fail, as does the test \code{test_one_way_road}. The planner answers $98$ seconds
from the Library to the Town Square, straight up the one-way road, the wrong way; the right answer is $204$ seconds,
through the Garden Centre.
\end{solution}

\begin{solution}{exo:g12:project-route-planner:7}
A \omterm{def:g10:graphs-first-look:graph}{graph} with the same arcs, each of cost $1$: the cost of a route is then its number of roads. From the Old Mill to
the Post Office it gives $7$ roads, through the Football Ground, the Swimming Pool, North Bridge, the Town Square, the
Garden Centre and the East Gate. That route is $8420$ m long and takes $770$ s; the fastest route has $10$ roads,
$9130$ m and $675$ s. Fewer roads, and even fewer metres, but over a minute and a half slower. Breadth-first search
would give the same answer, since all costs are equal; the tests compute both numbers with \code{cost_of}.
\end{solution}

\begin{solution}{exo:g12:project-route-planner:8}
Every other cost is in seconds, so the planner reads \code{0.25} as a quarter of a second: the shuttle looks almost
free, and every route that can use it does, for routes that in truth take fifteen minutes longer. Python adds the
numbers without complaint; nothing but the programmer knows the unit. The cost must be in the \omterm{def:g10:graphs-first-look:graph}{graph}'s unit:
$0.25 \times 3600 = 900$ seconds. Naming the unit in the code, as in \code{length_m} or \code{travel_seconds}, makes
such a mix visible.
\end{solution}

\begin{solution}{exo:g12:project-route-planner:9}
Each run finishes the $n$ vertices one by one, and each choice scans at most $n$ tentative costs: at most $n^2$
comparisons, plus one test per arc. For $n = 30$, about $900$ per run, so about $30 \times 900 = 27\,000$ for the
measure, plus $30 \times 105$ arc tests. Building a \omterm{def:g10:graphs-first-look:graph}{graph} reads $30$ places
and $53$ roads, a cost of the same order as one run. Building it once would save little here; on a \omterm{def:g10:graphs-first-look:graph}{graph} of millions
of arcs, it would matter.
\end{solution}

\begin{solution}{exo:g12:project-route-planner:10}
\smallskip\noindent\hspace*{2em}\code{def roads_on_route(conn, ids):}

\noindent\hspace*{4em}\code{found = []}

\noindent\hspace*{4em}\code{for a, b in zip(ids, ids[1:]):}

\noindent\hspace*{6em}\code{(road,) = conn.execute("SELECT road_id FROM road WHERE (from_place,"}

\noindent\hspace*{10em}\code{" to_place) IN ((?, ?), (?, ?))", (a, b, b, a)).fetchone()}

\noindent\hspace*{6em}\code{found.append(road)}

\noindent\hspace*{4em}\code{return found}

\smallskip\noindent The route from the Station to the Library uses roads $6$, $30$, $2$ and $1$. The tests close each
of the $49$ others in turn: the time stays $276$ s. The argument: closing a road only removes arcs, so no new route
appears and the best time cannot drop. And the old fastest route uses none of the removed arcs, so it is still there,
with the same cost: the best time cannot rise either.
\end{solution}

\begin{solution}{exo:g12:project-route-planner:11}
After each prompt the menu waits for one line; the run prints

\smallskip\noindent\hspace*{2em}\texttt{1 min 38 s over 1 road}

\noindent\hspace*{2em}\texttt{\ \ Town Square > Library}

\noindent\hspace*{2em}\texttt{3 min 24 s over 2 roads}

\noindent\hspace*{2em}\texttt{\ \ Library > Garden Centre > Town Square}

\smallskip\noindent each answer after its prompts, and the program ends at \code{4}. The way there is the one-way
road, $98$ s; the way back cannot use it.
\end{solution}

\begin{solution}{exo:g12:project-route-planner:12}
Build the \omterm{def:g10:graphs-first-look:graph}{graph} as usual, then remove the place: delete it from the successors of every vertex, and delete its own
entry. \omterm{met:g12:graph-algorithms:dijkstra}{Dijkstra's algorithm} on what remains gives the answer (\code{route_avoiding} in \code{exercises/rp_ex.py}).
Avoiding North Bridge, the fastest route from the Station to the Library goes round by the east: River Park, the Town
Hall, the Sports Hall, the East Gate and the Garden Centre, $572$ s instead of $276$.
\end{solution}

\begin{solution}{exo:g12:project-route-planner:13}
\emph{A route.} Each vertex but the source receives a parent $u$ only while the arc $u \to v$ is being examined, so
parent and child are joined by an arc; following parents from $b$ reaches the source, whose parent is \code{None},
since each parent was finished before its child, so the chain never loops. Reversed, the list is a route from $a$ to $b$.

\emph{Its cost.} When $v$ gets parent $u$, its tentative cost is set to $d(u) + c(u, v)$, where $d(u)$ is final, $u$
being finished. If $v$ later gets a better parent, both change together. So at the end $d(v) = d(\text{parent}) +
c(\text{parent}, v)$ for each $v$ on the route, and adding these equalities along the route gives $d(b)$ = the sum of
the costs of its arcs.

\emph{The fastest.} The costs are times, all $\ge 0$, so the theorem of chapter \ref{ch:g12:graph-algorithms} applies:
the cost of each finished vertex is its distance from the source, the smallest cost of any route. Hence no route from
$a$ to $b$ costs less than $d(b)$, the cost of the route returned.
\end{solution}

\begin{solution}{exo:g12:project-route-planner:14}
\smallskip\noindent\hspace*{2em}\code{def within(planner, start, seconds):}

\noindent\hspace*{4em}\code{dist, _ = dijkstra(from_database(planner.conn), planner.place_id(start))}

\noindent\hspace*{4em}\code{return sorted(planner.names[p] for p, t in dist.items()}

\noindent\hspace*{11em}\code{if 0 < t <= seconds)}

\smallskip\noindent Within $120$ s of the Station: Hill Park, Market Hall, North Bridge, River Park and the Town Hall.
The Town Hall is reached in $68$ s through River Park, not in $120$ s by its direct road. One run gives the time to
every place, so any limit is a filter on its result. (The run could also stop as soon as the next place to finish is
beyond the limit, since the places are finished in order of time.)
\end{solution}

\begin{solution}{exo:g12:project-route-planner:15}
The Old Mill lies in the south-west corner of the map and the Post Office near the north-east edge. The fastest route
does not go straight: it climbs north to the Station, where two roads at $70$ km/h meet, then turns east, $9130$ m
of road for $4.9$ km in a straight line. Its end is slow too: all three roads of the Post
Office are at $30$ km/h, so its last $680$ m alone take $82$ s. The South School and the Garden Centre are farther
apart, but their route keeps to roads at $50$ and $70$ km/h all the way. Time on a road map depends on
the roads, not on the distance between the ends.
\end{solution}
